% Unit 043 English translation of chapter3.tex lines 2215--2552.
% Source: Methods of Algebra, Volume 2, commit
% 9a5803ff77dd3257484cb177f851a73770a59dd3 (CC BY 4.0).
% stable-unit-id: o014.aljabr2.chapter3.classical-derived-functors
% Coverage: Section 3.12 in full.

% segment-id: o014.li2.chapter3.classical-derived-functors.h001
\section{Classical derived functors}\label{sec:derived-primer}
\hypertarget{o014.aljabr2.chapter3.classical-derived-functors}{ }

% segment-id: o014.li2.chapter3.classical-derived-functors.p001
The purpose of this section is to explain, without using the language of
derived categories, how to study right derived functors (respectively left
derived functors) for abelian categories with enough injectives (respectively
projectives). This already encompasses many cohomology theories commonly
encountered in applications. The definitions use the injective resolutions
(respectively projective resolutions) introduced in \S\ref{sec:resolutions}.

% segment-id: o014.li2.chapter3.classical-derived-functors.de001
\begin{definition}\label{def:derived-primer}
	\index{derived functor}
	\index[sym1]{RnF@$\mathrm{R}^n F$}
	\index[sym1]{LnF@$\mathrm{L}_n F$}
	Let $F:\mathcal{A}\to\mathcal{B}$ be an additive functor between abelian
	categories.
% segment-id: o014.li2.chapter3.classical-derived-functors.l001
	\begin{itemize}
% segment-id: o014.li2.chapter3.classical-derived-functors.i001
		\item Suppose that $\mathcal{A}$ has enough injectives. For every
		$X\in\cate{C}^+(\mathcal{A})$, choose an injective resolution
		$X\to I$. For each $n\in\Z$, define the value at $X$ of the
		\emph{$n$-th right derived functor} $\mathrm{R}^nF$ by
% segment-id: o014.li2.chapter3.classical-derived-functors.q001
		\[ \mathrm{R}^n F(X) := \Hm^n(\cate{C}F(I)). \]
% segment-id: o014.li2.chapter3.classical-derived-functors.i002
		\item Suppose that $\mathcal{A}$ has enough projectives. For every
		$X\in\cate{C}^-(\mathcal{A})$, choose a projective resolution
		$P\to X$. For each $n\in\Z$, define the value at $X$ of the
		\emph{$n$-th left derived functor} $\mathrm{L}_nF$ by
% segment-id: o014.li2.chapter3.classical-derived-functors.q002
		\[ \mathrm{L}_n F(X) := \Hm^{-n}(\cate{C}F(P)). \]
	\end{itemize}
	These functors all take values in $\mathcal{B}$. By construction, they come
	with canonical families of morphisms
	$\Hm^n(\cate{C}F(X))\to\mathrm{R}^nF(X)$ and
	$\mathrm{L}_nF(X)\to\Hm^{-n}(\cate{C}F(X))$.
\end{definition}

% segment-id: o014.li2.chapter3.classical-derived-functors.p002
This formulation leaves several questions. In what sense do these objects
depend on the chosen resolutions? And how does one promote a definition on
objects to a functor? Right and left derived functors are plainly dual, so we
discuss only right derived functors below. We begin with the second question,
encapsulating the basic tool in a lemma.

% segment-id: o014.li2.chapter3.classical-derived-functors.le001
\begin{lemma}\label{prop:resolution-homotopy-classical}
	Consider the solid part of the following diagram in
	$\cate{C}^+(\mathcal{A})$:
% segment-id: o014.li2.chapter3.classical-derived-functors.g001
	\[\begin{tikzcd}
		X \arrow[d, "f"'] \arrow[r] & I \arrow[dashed, d, "\beta"]\\
		Y \arrow[r] & J
	\end{tikzcd} \qquad \begin{array}{rl}
		X \to I: & \text{a quasi-isomorphism}, \\
		Y \to J: & \text{an injective resolution}, \\
		f: & \text{an arbitrary morphism}.
	\end{array}\]
	There is a unique morphism $\beta$ in $\cate{K}(\mathcal{A})$ for which
	the diagram commutes in $\cate{K}(\mathcal{A})$.
\end{lemma}
% segment-id: o014.li2.chapter3.classical-derived-functors.pf001
\begin{proof}
	Apply Theorem~\ref{prop:resolution-homotopy} to $X\to I$ and to the
	composite $X\xrightarrow{f}Y\to J$.
\end{proof}

% segment-id: o014.li2.chapter3.classical-derived-functors.p003
Reversing arrows gives the projective-resolution version in
$\cate{C}^-(\mathcal{A})$. In fact, the exercises for this chapter show that
when $X\hookrightarrow I$, one may choose $\beta$ so that the diagram commutes
in $\cate{C}(\mathcal{A})$. Since this section uses only $\Hm^n(\beta)$, the
version in $\cate{K}(\mathcal{A})$ suffices.

% segment-id: o014.li2.chapter3.classical-derived-functors.p004
Return to derived functors. Let $f:X\to Y$ be a morphism in
$\cate{C}^+(\mathcal{A})$, and choose injective resolutions $X\to I$ and
$Y\to J$. Lemma~\ref{prop:resolution-homotopy-classical} shows that, for every
$n\in\Z$, the morphism
% segment-id: o014.li2.chapter3.classical-derived-functors.q003
\[ \mathrm{R}^n F(f) := \Hm^n(\cate{K}F\beta): \; \Hm^n(\cate{K}F(I)) \to \Hm^n(\cate{K}F(J)) \]
depends only on $f$ and the resolutions $X\to I$, $Y\to J$. Once the
resolutions are fixed, one has
% segment-id: o014.li2.chapter3.classical-derived-functors.q004
\begin{gather*}
	\mathrm{R}^n F(f_1 + f_2) = \mathrm{R}^n F(f_1) + \mathrm{R}^n F(f_2), \\
	\mathrm{R}^n F(gf) = \mathrm{R}^n F(g) \; \mathrm{R}^n F(f), \quad \mathrm{R}^n F(\identity) = \identity.
\end{gather*}

% segment-id: o014.li2.chapter3.classical-derived-functors.p005
Taking $X=Y$ and $f=\identity_X$ also shows that different injective
resolutions give the same $(\mathrm{R}^nF)(X)$ up to a unique isomorphism.
Thus, like objects characterized by universal properties throughout category
theory, $\mathrm{R}^nF(X)$ is independent of its auxiliary data (the injective
resolution) up to canonical isomorphism, and
$\mathrm{R}^nF:\cate{C}^+(\mathcal{A})\to\mathcal{B}$ is an additive functor.

% segment-id: o014.li2.chapter3.classical-derived-functors.p006
By construction, if $F,G:\mathcal{A}\to\mathcal{B}$ are additive functors,
every morphism $F\to G$ canonically induces
$\mathrm{R}^nF\to\mathrm{R}^nG$ for all $n\in\Z$. Once an injective resolution
$X\to I$ is fixed, the morphism
$\mathrm{R}^nF(X)\to\mathrm{R}^nG(X)$ is concretely determined by
$\cate{C}F(I)\to\cate{C}G(I)$. Similarly, it induces
$\mathrm{L}_nF\to\mathrm{L}_nG$.

% segment-id: o014.li2.chapter3.classical-derived-functors.cn001
\begin{convention}
	Whenever we speak below of right derived functors (respectively left
	derived functors) of an additive functor
	$F:\mathcal{A}\to\mathcal{B}$, we implicitly assume that $\mathcal{A}$ has
	enough injectives (respectively enough projectives).
\end{convention}

% segment-id: o014.li2.chapter3.classical-derived-functors.th001
\begin{theorem}[Long exact sequence of derived functors]\label{prop:long-exact-sequence-primer}
	\index{long exact sequence!of derived functors}
	Let $F:\mathcal{A}\to\mathcal{B}$ be an additive functor between abelian
	categories, and consider a short exact sequence
	$0\to X\to Y\to Z\to0$ in $\cate{C}(\mathcal{A})$.
% segment-id: o014.li2.chapter3.classical-derived-functors.l002
	\begin{itemize}
% segment-id: o014.li2.chapter3.classical-derived-functors.i003
		\item If $X,Y,Z\in\Obj(\cate{C}^+(\mathcal{A}))$, there are canonical
		morphisms
		$\delta^n=\delta^n_{X,Y,Z}:\mathrm{R}^nF(Z)\to
		\mathrm{R}^{n+1}F(X)$ for $n\in\Z$, and an exact sequence
% segment-id: o014.li2.chapter3.classical-derived-functors.q005
		\[ \cdots \to \mathrm{R}^{n-1} F(Z) \xrightarrow{\delta^{n-1}} \mathrm{R}^n F(X) \to \mathrm{R}^n F(Y) \to \mathrm{R}^n F(Z) \xrightarrow{\delta^n} \mathrm{R}^{n+1} F(X) \to \cdots . \]
% segment-id: o014.li2.chapter3.classical-derived-functors.i004
		\item If $X,Y,Z\in\Obj(\cate{C}^-(\mathcal{A}))$, there are canonical
		morphisms
		$\partial_n=\partial_n^{X,Y,Z}:\mathrm{L}_nF(Z)\to
		\mathrm{L}_{n-1}F(X)$ for $n\in\Z$, and an exact sequence
% segment-id: o014.li2.chapter3.classical-derived-functors.q006
		\[ \cdots \to \mathrm{L}_{n+1} F(Z) \xrightarrow{\partial_{n+1}} \mathrm{L}_n F(X) \to \mathrm{L}_n F(Y) \to \mathrm{L}_n F(Z) \xrightarrow{\partial_n} \mathrm{L}_{n-1} F(X) \to \cdots . \]
	\end{itemize}
	The connecting morphisms $\delta^n$ and $\partial_n$ are functorial as
	follows. If
% segment-id: o014.li2.chapter3.classical-derived-functors.g002
	\[\begin{tikzcd}
		0 \arrow[r] & X \arrow[d, "\alpha"] \arrow[r] & Y \arrow[d, "\beta"] \arrow[r] & Z \arrow[d, "\gamma"] \arrow[r] & 0 \\
		0 \arrow[r] & X' \arrow[r] & Y' \arrow[r] & Z' \arrow[r] & 0
	\end{tikzcd}\]
	is a commutative diagram with exact rows, then, in the respective cases,
	the two diagrams
% segment-id: o014.li2.chapter3.classical-derived-functors.g003
	\[\begin{tikzcd}
		\mathrm{R}^n F(Z) \arrow[r, "{\delta^n}"] \arrow[d, "{\mathrm{R}^n F(\gamma)}"'] & \mathrm{R}^{n+1} F(X) \arrow[d, "{\mathrm{R}^{n+1} F(\alpha)}"] \\
		\mathrm{R}^n F(Z') \arrow[r, "{\delta^n}"'] & \mathrm{R}^{n+1} F(X')
	\end{tikzcd} \quad \text{and} \quad
% segment-id: o014.li2.chapter3.classical-derived-functors.g004
	\begin{tikzcd}
		\mathrm{L}_n F(Z) \arrow[r, "{\partial_n}"] \arrow[d, "{\mathrm{L}_n F(\gamma)}"'] & \mathrm{L}_{n-1} F(X) \arrow[d, "{\mathrm{L}_{n-1} F(\alpha)}"] \\
		\mathrm{L}_n F(Z') \arrow[r, "{\partial_n}"'] & \mathrm{L}_{n-1} F(X')
	\end{tikzcd}\]
	commute for every $n\in\Z$.
\end{theorem}
% segment-id: o014.li2.chapter3.classical-derived-functors.pf002
\begin{proof}
	We treat only $\mathrm{R}^nF$. Apply Proposition~\ref{prop:horseshoe} to
	$0\to X\to Y\to Z\to0$ to obtain a commutative diagram with exact rows
% segment-id: o014.li2.chapter3.classical-derived-functors.g005
	\[\begin{tikzcd}
		0 \arrow[r] & I \arrow[r] & J \arrow[r] & K \arrow[r] & 0 \\
		0 \arrow[r] & X \arrow[r] \arrow[u] & Y \arrow[r] \arrow[u] & Z \arrow[r] \arrow[u] & 0
	\end{tikzcd}\]
	in which every column is an injective resolution.
	Lemma~\ref{prop:inj-proj-split} ensures that
	$0\to I^n\to J^n\to K^n\to0$ splits. Its image under $F$ is therefore
	still a split short exact sequence, so
	$0\to\cate{C}F(I)\to\cate{C}F(J)\to\cate{C}F(K)\to0$ remains exact.
	Proposition~\ref{prop:long-exact-sequence-ses} now gives the morphisms
	$\delta^n$ and the long exact sequence
% segment-id: o014.li2.chapter3.classical-derived-functors.q007
	\[ \cdots \to \mathrm{R}^{n-1} F(Z) \xrightarrow{\delta^{n-1}} \mathrm{R}^n F(X) \to \mathrm{R}^n F(Y) \to \mathrm{R}^n F(Z) \xrightarrow{\delta^n} \cdots \]

	Functoriality of $\delta^n$ is more involved. We use the abelian category
	$\mathcal{A}^{\mathbf{2}}$ of Example~\ref{eg:A2-injective-projective}.
	First express the given commutative diagram with exact rows as the short
	exact sequence in $\mathcal{A}^{\mathbf{2}}$
% segment-id: o014.li2.chapter3.classical-derived-functors.q008
	\[ 0 \to [X \xrightarrow{\alpha} X'] \to [Y \xrightarrow{\beta} Y'] \to [Z \xrightarrow{\gamma} Z'] \to 0. \]
	Example~\ref{eg:A2-injective-projective} says that
	$\mathcal{A}^{\mathbf{2}}$ has enough injectives. Choose injective
	resolutions $[X\to X']\hookrightarrow\mathcal{I}$ and
	$[Z\to Z']\hookrightarrow\mathcal{K}$, and use
	Proposition~\ref{prop:horseshoe} to extend them to a commutative diagram
	with exact rows in $\cate{C}(\mathcal{A}^{\mathbf{2}})$:
% segment-id: o014.li2.chapter3.classical-derived-functors.g006
	\begin{equation}\label{eqn:long-exact-sequence-primer-aux}\begin{tikzcd}
		0 \arrow[r] & \mathcal{I} \arrow[r] & \mathcal{J} \arrow[r] & \mathcal{K} \arrow[r] & 0 \\
		0 \arrow[r] & {[X \to X']} \arrow[r] \arrow[hookrightarrow, u] & {[Y \to Y']} \arrow[r] \arrow[hookrightarrow, u] & {[Z \to Z']} \arrow[r] \arrow[hookrightarrow, u] & 0
	\end{tikzcd}\end{equation}
	such that $[Y\to Y']\to\mathcal{J}$ is also an injective resolution.
	Write $\mathcal{J}^n=[J^n\to(J')^n]$, and similarly for the other terms.
	The evaluation functors $\mathrm{ev}_0$ and $\mathrm{ev}_1$ from
	Example~\ref{eg:A2-injective-projective} show that $J^n$, $(J')^n$, and
	the corresponding terms are injective objects of $\mathcal{A}$, while
	$Y\to J$, $Y'\to J'$, and the corresponding maps are injective
	resolutions in $\cate{C}(\mathcal{A})$. Thus compatible injective
	resolutions have been erected above the original diagram.

	Now expand the first row of
	\eqref{eqn:long-exact-sequence-primer-aux} as the following commutative
	diagram with exact rows in $\cate{C}(\mathcal{A})$:
% segment-id: o014.li2.chapter3.classical-derived-functors.g007
	\[\begin{tikzcd}
		0 \arrow[r] & I \arrow[d] \arrow[r] & J \arrow[d] \arrow[r] & K \arrow[d] \arrow[r] & 0 \\
		0 \arrow[r] & I' \arrow[r] & J' \arrow[r] & K' \arrow[r] & 0
	\end{tikzcd}\]
	For every $n$, both rows split because $I^n$ and $(I')^n$ are injective.
	Applying $F$ degreewise therefore preserves their exactness. Functoriality
	of $\delta^n$ now reduces to the corresponding assertion of
	Proposition~\ref{prop:long-exact-sequence-ses}.
\end{proof}

% segment-id: o014.li2.chapter3.classical-derived-functors.p007
In the classical setting, one is mainly interested in the restrictions of
$\mathrm{R}^nF$ and $\mathrm{L}_nF$ to $\mathcal{A}$, still written
$\mathrm{R}^nF(X)$ and $\mathrm{L}_nF(X)$ for
$X\in\Obj(\mathcal{A})$. They are computed by flattening an injective
resolution to an exact sequence $0\to X\to I^0\to I^1\to\cdots$ (or a
projective resolution to $\cdots\to P^1\to P^0\to X\to0$); see
Example~\ref{eg:resolution-single}. Derived functors defined on complexes were
called \emph{hyper-derived functors} in the early literature and are more
naturally treated using derived categories or spectral sequences. We therefore
focus below on the classical case $X\in\Obj(\mathcal{A})$.
\index{hyper-derived functor}

% segment-id: o014.li2.chapter3.classical-derived-functors.p008
The first step is to distill the long exact sequence of
Theorem~\ref{prop:long-exact-sequence-primer} into the notion of a
$\delta$-functor.

% segment-id: o014.li2.chapter3.classical-derived-functors.de002
\begin{definition}\label{def:delta-functor}
	\index{$\delta$-functor}
	Let $\mathcal{A}$ and $\mathcal{B}$ be abelian categories. A
	\emph{cohomological $\delta$-functor} from $\mathcal{A}$ to $\mathcal{B}$
	consists of the following data:
% segment-id: o014.li2.chapter3.classical-derived-functors.l003
	\begin{compactitem}
% segment-id: o014.li2.chapter3.classical-derived-functors.i005
		\item additive functors $F^n:\mathcal{A}\to\mathcal{B}$ for
		$n\in\Z_{\geq0}$;
% segment-id: o014.li2.chapter3.classical-derived-functors.i006
		\item for each short exact sequence $0\to X\to Y\to Z\to0$ in
		$\mathcal{A}$, specified morphisms $\delta^n:F^nZ\to F^{n+1}X$,
		called connecting morphisms, for $n\in\Z_{\geq0}$. They are
		functorial in morphisms of short exact sequences and yield an exact
		sequence
% segment-id: o014.li2.chapter3.classical-derived-functors.g008
		\begin{equation*}\begin{tikzcd}[row sep=large]
			0 \arrow[r] & F^0(X) \arrow[r] & F^0(Y) \arrow[r] \arrow[phantom, d, ""{coordinate, name=A}] & F^0(Z) \arrow[dll, rounded corners, "{\delta^0}" description, to path={
				-- ([xshift=2ex]\tikztostart.east)
				|- (A) [near end]\tikztonodes
				-| ([xshift=-2ex]\tikztotarget.west)
				-- (\tikztotarget)}] \\
			& F^1(X) \arrow[r] & F^1(Y) \arrow[r] \arrow[phantom, d, ""{coordinate, name=B}] & F^1(Z) \arrow[dll, rounded corners, "{\delta^1}" description, to path={
				-- ([xshift=2ex]\tikztostart.east)
				|- (B) [near end]\tikztonodes
				-| ([xshift=-2ex]\tikztotarget.west)
				-- (\tikztotarget)}] \\
			& F^2(X) \arrow[r] & F^2(Y) \arrow[r] & \cdots .
		\end{tikzcd}\end{equation*}
	\end{compactitem}
	A morphism from a cohomological $\delta$-functor
	$(F^n,\delta^n)_{n\geq0}$ to $(G^n,\eta^n)_{n\geq0}$ is a family of
	morphisms $\varphi^n:F^n\to G^n$ compatible with the connecting morphisms
	of short exact sequences. In other words, the diagram
% segment-id: o014.li2.chapter3.classical-derived-functors.g009
	\[\begin{tikzcd}
		G^{n-1} (Z) \arrow[r, "{\eta^{n-1}}"] & G^n (X) \\
		F^{n-1} (Z) \arrow[r, "{\delta^{n-1}}"'] \arrow[u, "{\varphi^{n-1}}"] & F^n (X) \arrow[u, "\varphi^n"']
	\end{tikzcd}\]%
	\footnote{Translator's note (O014-C060): both objects in the left column
	are printed as $X$ in the source. For a short exact sequence
	$0\to X\to Y\to Z\to0$, the connecting morphisms $\delta^{n-1}$ and
	$\eta^{n-1}$ have domains $F^{n-1}(Z)$ and $G^{n-1}(Z)$, respectively,
	and codomains $F^n(X)$ and $G^n(X)$. Both left-hand objects have therefore
	been restored to $Z$.}
	must commute for all $n\geq1$ and every short exact sequence
	$0\to X\to Y\to Z\to0$.

	Dually, a \emph{homological $\delta$-functor} from $\mathcal{A}$ to
	$\mathcal{B}$ consists of the following data:
% segment-id: o014.li2.chapter3.classical-derived-functors.l004
	\begin{compactitem}
% segment-id: o014.li2.chapter3.classical-derived-functors.i007
		\item additive functors $F_n:\mathcal{A}\to\mathcal{B}$ for
		$n\in\Z_{\geq0}$;
% segment-id: o014.li2.chapter3.classical-derived-functors.i008
		\item for each short exact sequence $0\to X\to Y\to Z\to0$ in
		$\mathcal{A}$, specified connecting morphisms
		$\partial_n:F_nZ\to F_{n-1}X$ for $n\geq1$, functorial in morphisms of
		short exact sequences and yielding an exact sequence
% segment-id: o014.li2.chapter3.classical-derived-functors.g010
		\begin{equation*}\begin{tikzcd}[row sep=large]
			\cdots \arrow[r] & F_2(Y) \arrow[r] \arrow[phantom, d, ""{coordinate, name=A}] & F_2(Z) \arrow[dll, rounded corners, "{\partial_2}" description, to path={
				-- ([xshift=2ex]\tikztostart.east)
				|- (A) [near end]\tikztonodes
				-| ([xshift=-2ex]\tikztotarget.west)
				-- (\tikztotarget)}] & \\
			F_1(X) \arrow[r] & F_1(Y) \arrow[r] \arrow[phantom, d, ""{coordinate, name=B}] & F_1(Z) \arrow[dll, rounded corners, "{\partial_1}" description, to path={
				-- ([xshift=2ex]\tikztostart.east)
				|- (B) [near end]\tikztonodes
				-| ([xshift=-2ex]\tikztotarget.west)
				-- (\tikztotarget)}] & \\
			F_0 (X) \arrow[r] & F_0 (Y) \arrow[r] & F_0 (Z) \arrow[r] & 0.
		\end{tikzcd}\end{equation*}
	\end{compactitem}
	Morphisms of homological $\delta$-functors are defined analogously, with
	compatibility with connecting morphisms.
\end{definition}

% segment-id: o014.li2.chapter3.classical-derived-functors.p009
The long exact sequences in the definition automatically imply that $F^0$ is
left exact and $F_0$ is right exact. We now return to derived functors.

% segment-id: o014.li2.chapter3.classical-derived-functors.pr001
\begin{proposition}\label{prop:derived-primer-long}
	Let $F:\mathcal{A}\to\mathcal{B}$ be an additive functor between abelian
	categories. If $\mathcal{A}$ has enough injectives, then:
% segment-id: o014.li2.chapter3.classical-derived-functors.l005
	\begin{compactitem}
% segment-id: o014.li2.chapter3.classical-derived-functors.i009
		\item $n<0\implies\mathrm{R}^nF=0$;
% segment-id: o014.li2.chapter3.classical-derived-functors.i010
		\item if $I$ is injective in $\mathcal{A}$, then
		$n>0\implies\mathrm{R}^nF(I)=0$;
% segment-id: o014.li2.chapter3.classical-derived-functors.i011
		\item $(\mathrm{R}^nF,\delta^n)_{n\geq0}$ is a cohomological
		$\delta$-functor;
% segment-id: o014.li2.chapter3.classical-derived-functors.i012
		\item if $F$ is left exact, there is a canonical isomorphism
		$F\rightiso\mathrm{R}^0F$.
	\end{compactitem}
	Dually, if $\mathcal{A}$ has enough projectives, then:
% segment-id: o014.li2.chapter3.classical-derived-functors.l006
	\begin{compactitem}
% segment-id: o014.li2.chapter3.classical-derived-functors.i013
		\item $n<0\implies\mathrm{L}_nF=0$;
% segment-id: o014.li2.chapter3.classical-derived-functors.i014
		\item if $P$ is projective in $\mathcal{A}$, then
		$n>0\implies\mathrm{L}_nF(P)=0$;
% segment-id: o014.li2.chapter3.classical-derived-functors.i015
		\item $(\mathrm{L}_nF,\partial_n)_{n\geq0}$ is a homological
		$\delta$-functor;
% segment-id: o014.li2.chapter3.classical-derived-functors.i016
		\item if $F$ is right exact, there is a canonical isomorphism
		$\mathrm{L}_0F\rightiso F$.
	\end{compactitem}
\end{proposition}
% segment-id: o014.li2.chapter3.classical-derived-functors.pf003
\begin{proof}
	By duality, it suffices to treat $\mathrm{R}^nF$. Choose an injective
	resolution $0\to X\to I^0\to I^1\to\cdots$ of
	$X\in\Obj(\mathcal{A})$. Then
% segment-id: o014.li2.chapter3.classical-derived-functors.q009
	\[ \mathrm{R}^n F(X) = \Hm^n ( \cdots \to 0 \to 0 \to \underbracket{FI^0}_{\text{term in degree $0$}} \to \underbracket{FI^1}_{\text{term in degree $1$}} \to \cdots ). \]
	This immediately gives $\mathrm{R}^nF(X)=0$ when $n<0$. If $X=I$ is
	injective, take the resolution $0\to I\xrightarrow{\identity}I\to0\to
	\cdots$, which gives $\mathrm{R}^nF(I)=0$ for $n>0$.

	In view of this, the assertion about cohomological $\delta$-functors is
	just a restatement of Theorem~\ref{prop:long-exact-sequence-primer}.

	Finally suppose that $F$ is left exact. Since
	$X\rightiso\Ker[I^0\to I^1]$ and a left exact functor preserves kernels,
	$\mathrm{R}^0F(X)=\Ker[FI^0\to FI^1]$ is canonically isomorphic to $FX$.
\end{proof}

% segment-id: o014.li2.chapter3.classical-derived-functors.p010
In general, one considers right derived functors only for left exact functors,
and left derived functors only for right exact functors. Why? An elementary
explanation comes from the problem of extending exact sequences. For example,
if $F$ is left exact and $0\to X\to Y\to Z\to0$ is a short exact sequence in
$\mathcal{A}$, we would like to extend the exact sequence
$0\to FX\to FY\to FZ$ as far as possible to the right. Right derived
functors do exactly this.

% segment-id: o014.li2.chapter3.classical-derived-functors.p011
For $n\geq1$, the functors $\mathrm{R}^nF$ and $\mathrm{L}_nF$ are
customarily called the \emph{higher derived functors} of $F$, with domain
understood to be the objects of $\mathcal{A}$ rather than complexes. As an
application of extending an exact sequence, the following corollary says that
higher derived functors are precisely the obstructions to exactness.

% segment-id: o014.li2.chapter3.classical-derived-functors.co001
\begin{corollary}\label{prop:obstruction-exactness}
	For a left exact (respectively right exact) additive functor
	$F:\mathcal{A}\to\mathcal{B}$, the following are equivalent.
% segment-id: o014.li2.chapter3.classical-derived-functors.l007
	\begin{enumerate}[(i)]
% segment-id: o014.li2.chapter3.classical-derived-functors.i017
		\item $F:\mathcal{A}\to\mathcal{B}$ is exact.
% segment-id: o014.li2.chapter3.classical-derived-functors.i018
		\item For every $X\in\Obj(\mathcal{A})$ and $n>0$,
		$\mathrm{R}^nF(X)=0$ (respectively $\mathrm{L}_nF(X)=0$).
% segment-id: o014.li2.chapter3.classical-derived-functors.i019
		\item For every $X\in\Obj(\mathcal{A})$,
		$\mathrm{R}^1F(X)=0$ (respectively $\mathrm{L}_1F(X)=0$).
	\end{enumerate}
\end{corollary}
% segment-id: o014.li2.chapter3.classical-derived-functors.pf004
\begin{proof}
	Consider the left exact case. For (i) $\implies$ (ii), the complex
	$I^0\to I^1\to\cdots$ used to compute the derived functors is exact away
	from its zeroth term, and therefore so is its image under $F$. The
	implication (ii) $\implies$ (iii) is trivial. For (iii) $\implies$ (i), use
	$\mathrm{R}^0F\simeq F$ and apply the exact sequence
% segment-id: o014.li2.chapter3.classical-derived-functors.q010
	\[ 0 \to FX \to FY \to FZ \to \mathrm{R}^1 F(X) \]
	to each short exact sequence $0\to X\to Y\to Z\to0$ in $\mathcal{A}$.
\end{proof}

% segment-id: o014.li2.chapter3.classical-derived-functors.cn002
\begin{convention}\label{con:F-acyclic}\index{$F$-acyclic object}
	Let $F:\mathcal{A}\to\mathcal{B}$ be a left exact (respectively right
	exact) additive functor. An object $A\in\Obj(\mathcal{A})$ is called
	\emph{$F$-acyclic} if all its higher right derived functors (respectively
	left derived functors) under $F$ vanish. For example, when $F$ is left
	exact (respectively right exact), Proposition~\ref{prop:derived-primer-long}
	implies that every injective (respectively projective) object is
	$F$-acyclic.
\end{convention}

% segment-id: o014.li2.chapter3.classical-derived-functors.p012
The following classical technique, called dimension shifting, is often used in
recursive arguments with derived functors.

% segment-id: o014.li2.chapter3.classical-derived-functors.pr002
\begin{proposition}[Dimension shifting]\label{prop:dimension-shifting}
	\index{dimension shifting}
	Let $F:\mathcal{A}\to\mathcal{B}$ be a left exact (respectively right
	exact) additive functor, and suppose that $\mathcal{A}$ has enough
	injectives (respectively enough projectives). Let
% segment-id: o014.li2.chapter3.classical-derived-functors.q011
	\[ 0 \to X \to A \to B \to 0, \quad (\text{respectively}\; 0 \to B \to A \to X \to 0 ), \]
	be a short exact sequence in $\mathcal{A}$, with $A$ $F$-acyclic. Then for
	$n\geq1$ there are natural isomorphisms
% segment-id: o014.li2.chapter3.classical-derived-functors.q012
	\[\mathrm{R}^n F(X) \simeq \begin{cases}
		\mathrm{R}^{n-1} F(B), & n \geq 2, \\
		\Coker\left[FA \to FB \right], & n = 1,
	\end{cases}\]
	respectively
% segment-id: o014.li2.chapter3.classical-derived-functors.q013
	\[ \mathrm{L}_n F(X) \simeq \begin{cases}
		\mathrm{L}_{n-1} F(B), & n \geq 2, \\
		\Ker\left[ FB \to FA \right], & n = 1.
	\end{cases} \]
\end{proposition}
% segment-id: o014.li2.chapter3.classical-derived-functors.pf005
\begin{proof}
	By duality, prove only the assertion about $\mathrm{R}^nF(X)$. Inspect the
	corresponding long exact sequence
% segment-id: o014.li2.chapter3.classical-derived-functors.q014
	\[ \mathrm{R}^{n-1} F(A) \to \mathrm{R}^{n-1} F(B) \to \mathrm{R}^n F(X) \to \underbracket{\mathrm{R}^n F(A)}_{= 0}, \quad n \in \Z_{\geq 1}, \]
	and use $\mathrm{R}^0F(A)\simeq FA$ and $\mathrm{R}^0F(B)=FB$.
\end{proof}

% segment-id: o014.li2.chapter3.classical-derived-functors.p013
The next key result shows that $F$-acyclic objects may also be used to compute
derived functors.

% segment-id: o014.li2.chapter3.classical-derived-functors.co002
\begin{corollary}\label{prop:acyclic-resolution}
	Let $F:\mathcal{A}\to\mathcal{B}$ be left exact (respectively right
	exact), and suppose there is an exact sequence in $\mathcal{A}$
% segment-id: o014.li2.chapter3.classical-derived-functors.q015
	\[ 0 \to X  \to A^0 \to A^1 \to \cdots \quad (\text{respectively}\; \cdots \to A_1 \to A_0 \to X \to 0 ), \]
	where every $A^n$ (respectively $A_n$) is $F$-acyclic for $n\geq0$.
	In addition, set $A^{-n}=0$ (respectively $A_{-n}=0$) for every
	$n\geq1$.\footnote{Translator's note (O014-C062): the source gives no
	range for the convention $A^{-n}=0$ (respectively $A_{-n}=0$). If it were
	applied also at $n=0$, it would force $A^0=0$ (respectively $A_0=0$),
	contradicting the resolution term just introduced. The intended range
	$n\geq1$ has been made explicit.} Then, for every $n\geq0$, there are
	natural isomorphisms
% segment-id: o014.li2.chapter3.classical-derived-functors.q016
	\[ \mathrm{R}^n F(X) \simeq \Hm^n\left( FA^\bullet \right) \quad \text{or}\quad \mathrm{L}_n F(X) \simeq \Hm_n\left(FA_\bullet\right). \]
\end{corollary}
% segment-id: o014.li2.chapter3.classical-derived-functors.pf006
\begin{proof}
	We prove only the assertion about $\mathrm{R}^nF(X)$. For $m\geq1$, put
% segment-id: o014.li2.chapter3.classical-derived-functors.q017
	\[ B^m := \Image\left[ A^{m-1} \to A^m \right] = \Ker\left[ A^m \to A^{m+1} \right]; \]
	also put $B^0:=X$. The original exact sequence then decomposes into
% segment-id: o014.li2.chapter3.classical-derived-functors.q018
	\[ 0 \to B^m \to A^m \to B^{m+1} \to 0, \quad m \in \Z_{\geq 0}. \]
	When $n=0$, the claim reduces to preservation of kernels:
	$FX\simeq\Ker[FA^0\to FA^1]$. For $n\geq1$, repeated dimension shifting
	using Proposition~\ref{prop:dimension-shifting} gives natural isomorphisms
% segment-id: o014.li2.chapter3.classical-derived-functors.q019
	\[ \mathrm{R}^n F(X) = \mathrm{R}^n F(B^0) \simeq \cdots \simeq \mathrm{R}^1 F(B^{n-1}) \simeq \Coker\left[ FA^{n-1} \to FB^n \right]; \]
	and because $F$ preserves kernels, the definition of $B^n$ identifies
	$FB^n$ with $\Ker[FA^n\to FA^{n+1}]$. Thus, for $n\geq1$,
% segment-id: o014.li2.chapter3.classical-derived-functors.q020
	\[ \mathrm{R}^n F(X) \simeq \frac{\Ker\left[ FA^n \to FA^{n+1} \right]}{\Image\left[ FA^{n-1} \to FA^n \right]}. \]
\end{proof}

% segment-id: o014.li2.chapter3.classical-derived-functors.p014
The preceding argument relies mainly on the long exact sequence and hardly at
all on the definition of derived functors. This suggests characterizing the
derived functors of $F$ among general cohomological (respectively homological)
$\delta$-functors. The key requirement is that their positive-degree terms be
effaceable (respectively co-effaceable).

% segment-id: o014.li2.chapter3.classical-derived-functors.de003
\begin{definition}[A.\ Grothendieck]\label{def:erasable}
	\index{effaceable functor}
	\index{co-effaceable functor}
	Let $E:\mathcal{A}\to\mathcal{B}$ be an additive functor between abelian
	categories. It is called \emph{effaceable} (respectively
	\emph{co-effaceable}) if for every $X\in\Obj(\mathcal{A})$ there is a
	monomorphism $a:X\hookrightarrow Y$ (respectively an epimorphism
	$a:Y\twoheadrightarrow X$) such that $Ea=0$.
\end{definition}

% segment-id: o014.li2.chapter3.classical-derived-functors.p015
We shall derive the following universal property from effaceability.

% segment-id: o014.li2.chapter3.classical-derived-functors.de004
\begin{definition}\label{def:universal-delta-functor}
	\index{$\delta$-functor!universal}
	Let $(R^n,\delta^n)_{n\geq0}$ be a cohomological $\delta$-functor from
	$\mathcal{A}$ to $\mathcal{B}$. It is called a \emph{universal
	cohomological $\delta$-functor} if, for every cohomological $\delta$-functor
	$(F^n,\delta_F^n)_{n\geq0}$ from $\mathcal{A}$ to $\mathcal{B}$, every
	morphism $\varphi^0:R^0\to F^0$ extends uniquely to a morphism
	$(\varphi^n)_{n\geq0}$ of $\delta$-functors.

	Dually, let $(L_n,\partial_n)_{n\geq0}$ be a homological $\delta$-functor
	with the following property: for every homological $\delta$-functor
	$(F_n,\partial_n^F)_{n\geq0}$, every morphism $\varphi_0:F_0\to L_0$
	extends uniquely to $(\varphi_n)_{n\geq0}$. Then
	$(L_n,\partial_n)_{n\geq0}$ is called a \emph{universal homological
	$\delta$-functor}.
\end{definition}

% segment-id: o014.li2.chapter3.classical-derived-functors.p016
Our goal is to show that, for a left exact (respectively right exact) functor
$F:\mathcal{A}\to\mathcal{B}$, a universal cohomological $\delta$-functor
with $R^0=F$ (respectively a universal homological $\delta$-functor with
$L_0=F$), if it exists, is unique up to a unique isomorphism.

% segment-id: o014.li2.chapter3.classical-derived-functors.le002
\begin{lemma}\label{prop:erasable-aux}
	Let $E:\mathcal{A}\to\mathcal{B}$ be effaceable, let
	$0\to X\xrightarrow{u}I\xrightarrow{v}B\to0$ be a short exact sequence
	in $\mathcal{A}$, and let $f:X\to Y$ be any morphism. There is a
	commutative diagram with exact rows
% segment-id: o014.li2.chapter3.classical-derived-functors.g011
	\begin{equation*}\begin{tikzcd}
			0 \arrow[r] & X \arrow[d, "f"'] \arrow[r, "u"] & I \arrow[r, "v"] \arrow[d] & B \arrow[d] \arrow[r] & 0 \\
			0 \arrow[r] & Y \arrow[r, "\iota"'] & J \arrow[r] & C \arrow[r] & 0 ,
	\end{tikzcd}\end{equation*}
	such that $E\iota=0$.
\end{lemma}
% segment-id: o014.li2.chapter3.classical-derived-functors.pf007
\begin{proof}
	Choose a monomorphism $h:Y\hookrightarrow J_0$ such that $Eh=0$, and form
	the pushout diagram
% segment-id: o014.li2.chapter3.classical-derived-functors.g012
	\[\begin{tikzcd}
		X \arrow[hookrightarrow, r, "u"] \arrow[d, "hf"'] & I \arrow[d] \\
		J_0 \arrow[r, "k"'] & J \arrow[phantom, lu, "\boxplus" description] ,
	\end{tikzcd}\]
	where $k$ is a monomorphism by
	Proposition~\ref{prop:Abel-cat-pull-push}. Set $\iota=kh:Y\to J$ and
	$C:=\Coker(\iota)$. This gives the left-hand square of the required
	diagram; functoriality of cokernels gives the right-hand square.
\end{proof}

% segment-id: o014.li2.chapter3.classical-derived-functors.pr003
\begin{proposition}\label{prop:erasable-univ-delta}
	Let $(R^n,\delta^n)_{n\geq0}$ be a cohomological $\delta$-functor from
	$\mathcal{A}$ to $\mathcal{B}$. If $R^n$ is effaceable for every $n>0$,
	then $(R^n,\delta^n)_{n\geq0}$ is universal.

	Dually, a homological $\delta$-functor whose positive-degree parts are
	co-effaceable is universal.
\end{proposition}
% segment-id: o014.li2.chapter3.classical-derived-functors.pf008
\begin{proof}
	We prove only the cohomological version. Let
	$(F^n,\delta_F^n)_{n\geq0}$ be another cohomological $\delta$-functor from
	$\mathcal{A}$ to $\mathcal{B}$, and fix $\varphi^0:R^0\to F^0$. We
	construct $\varphi^n$ recursively for $n\geq1$. Given
	$X\in\Obj(\mathcal{A})$, the hypothesis supplies a short exact sequence
% segment-id: o014.li2.chapter3.classical-derived-functors.q021
	\[ 0 \to X \to I \to B \to 0 \]
	such that $R^n(X\to I)=0$. The corresponding long exact sequences give
	the solid part of a commutative diagram with exact rows
% segment-id: o014.li2.chapter3.classical-derived-functors.g013
	\[\begin{tikzcd}
		F^{n-1} I \arrow[r] & F^{n-1} B \arrow[r, "{\delta_F^{n-1}}"] & F^n X \arrow[r] & F^n I \\
		R^{n-1} I \arrow[u, "{\varphi^{n-1}_I}"] \arrow[r] & R^{n-1} B \arrow[u, "{\varphi^{n-1}_B}"] \arrow[r, "{\delta^{n-1}}"'] & R^n X \arrow[dashed, u, "{\varphi^n_X}"'] \arrow[r] & 0 .
	\end{tikzcd}\]
	Exactness of the rows and the universal property of the cokernel give a
	unique dashed arrow $\varphi_X^n$ making the whole diagram commute. This
	characterizes $\varphi_X^n$, although it still appears to depend on the
	choice of $X\hookrightarrow I$.

	Given any morphism $f:X\to Y$, we claim there is a commutative diagram
	with exact rows
% segment-id: o014.li2.chapter3.classical-derived-functors.g014
	\begin{equation}\label{eqn:erasable-univ-delta}\begin{tikzcd}
		0 \arrow[r] & X \arrow[d, "f"'] \arrow[r] & I \arrow[r] \arrow[d] & B \arrow[d] \arrow[r] & 0 \\
		0 \arrow[r] & Y \arrow[r] & J \arrow[r] & C \arrow[r] & 0 ,
	\end{tikzcd}\end{equation}
	such that both $R^n(X\to I)$ and $R^n(Y\to J)$ are zero. Indeed, first
	choose $X\hookrightarrow I$, then apply Lemma~\ref{prop:erasable-aux}.
	Substituting this into the constructions of $\varphi_X^n$ and
	$\varphi_Y^n$ yields
% segment-id: o014.li2.chapter3.classical-derived-functors.g015
	\[\begin{tikzcd}[row sep=small, column sep=small]
		& F^{n-1} B \arrow[rr] \arrow[dd] & & F^n X \arrow[dd] \\
		R^{n-1} B \arrow[ru] \arrow[rr, crossing over] \arrow[dd] & & R^n X \arrow[ru, "{\varphi_X^n}"'] & \\
		& F^{n-1} C \arrow[rr] & & F^n Y \\
		R^{n-1} C \arrow[ru] \arrow[rr] & & R^n Y \arrow[ru, "{\varphi_Y^n}"'] \arrow[leftarrow, uu, crossing over] &
	\end{tikzcd}\]
	Every face except the right one commutes by construction, definition, or
	the induction hypothesis. Since $R^{n-1}B\to R^nX$ is an epimorphism, the
	usual argument shows that the right face commutes too.

	This simultaneously proves that $\varphi_X^n$ is independent of the choice
	of $X\hookrightarrow I$ (take $f=\identity_X$ and observe from the proof of
	Lemma~\ref{prop:erasable-aux} that any two maps $X\hookrightarrow I_i$ can
	be sent by pushouts into a common $X\hookrightarrow J$) and that it is
	functorial in $X$ (take arbitrary $f$). This recursively completes the
	construction of $\varphi^n$.

	It remains to show that $(\varphi^n)_{n\geq0}$ is compatible with the
	connecting morphisms. Given a short exact sequence
	$0\to X\to Y\to Z\to0$, use Lemma~\ref{prop:erasable-aux} to choose a
	commutative diagram with exact rows
% segment-id: o014.li2.chapter3.classical-derived-functors.g016
	\[\begin{tikzcd}
		0 \arrow[r] & X \arrow[d, "\identity_X"'] \arrow[r] & Y \arrow[r] \arrow[d] & Z \arrow[d, "\alpha"] \arrow[r] & 0 \\
		0 \arrow[r] & X \arrow[r] & I \arrow[r] & B \arrow[r] & 0
	\end{tikzcd}\]
	such that $R^n(X\to I)=0$. For each $n\geq1$, form the diagram
% segment-id: o014.li2.chapter3.classical-derived-functors.g017
	\[\begin{tikzcd}[column sep=large]
		F^{n-1} Z \arrow[r, "{F^{n-1}\alpha}"] & F^{n-1} B \arrow[r, "{\delta_F^{n-1}}"] & F^n X \\
		R^{n-1} Z \arrow[u, "{\varphi_Z^{n-1}}"] \arrow[r, "{R^{n-1}\alpha}"'] & R^{n-1} B \arrow[u, "{\varphi_B^{n-1}}"] \arrow[r, "{\delta^{n-1}}"'] & R^n X \arrow[u, "{\varphi_X^n}"']
	\end{tikzcd}\]%
	\footnote{Translator's note (O014-C061): the source labels the two
	connecting arrows in this diagram $\delta_F^n$ and $\delta^n$. Since their
	domains are $F^{n-1}B$ and $R^{n-1}B$, and their codomains are $F^nX$ and
	$R^nX$, the well-typed indices are $n-1$; the labels have been restored to
	$\delta_F^{n-1}$ and $\delta^{n-1}$.}
	By the preceding commutative diagram and the axioms of a cohomological
	$\delta$-functor, the composites along the two rows are the connecting
	morphisms $F^{n-1}Z\to F^nX$ and $R^{n-1}Z\to R^nX$ whose compatibility
	is at issue. The left square commutes by functoriality of
	$\varphi^{n-1}$, and the right square by the construction of
	$\varphi_X^n$.
\end{proof}

% segment-id: o014.li2.chapter3.classical-derived-functors.co003
\begin{corollary}[Characterization of derived functors]\label{prop:derived-erasable}
	Let $F:\mathcal{A}\to\mathcal{B}$ be left exact and suppose that
	$\mathcal{A}$ has enough injectives. Then
	$(\mathrm{R}^nF,\delta^n)_{n\geq0}$ is the universal cohomological
	$\delta$-functor satisfying $\mathrm{R}^0F\simeq F$.

	Dually, if $F$ is right exact and $\mathcal{A}$ has enough projectives,
	then $(\mathrm{L}_nF,\partial_n)_{n\geq0}$ is the universal homological
	$\delta$-functor satisfying $\mathrm{L}_0F\simeq F$.
\end{corollary}
% segment-id: o014.li2.chapter3.classical-derived-functors.pf009
\begin{proof}
	It is enough to discuss the left exact case. For $n>0$, the functor
	$\mathrm{R}^nF$ is effaceable: every $X\in\Obj(\mathcal{A})$ embeds into
	an injective object $I$, and Proposition~\ref{prop:derived-primer-long}
	gives $\mathrm{R}^nF(I)=0$. The same proposition gives
	$\mathrm{R}^0F\simeq F$. Apply Proposition~\ref{prop:erasable-univ-delta}.
\end{proof}
